\chapter*{Заключение}
\addcontentsline{toc}{chapter}{Заключение}

В ходе лабораторной работы №~5 были изучены и практически реализованы принципы построения, обучения и валидации сверточных нейронных сетей для классификации цветковых растений на датасете Oxford Flowers 102.

\textbf{Основные результаты работы:}
\begin{enumerate}[itemsep=1pt, parsep=0pt]
    \item \textbf{Конвейер данных:} Реализован конвейер \texttt{tf.data} со стандартизацией $96 \times 96 \times 3$, масштабированием $[0.0, 1.0]$ и потоковой аугментацией (отражения, повороты и масштабирование).
    \item \textbf{Архитектура модели:} Построена сверточная нейросеть из 3 блоков свертки (32, 64, 128 фильтров $3 \times 3$), слоев BatchNormalization, нелинейностей ReLU и MaxPooling2D, объединенных с 2 полносвязными слоями (Dense на 256 нейронов и выходной классификатор на 102 класса с Softmax). Число параметров составило $4\,839\,750$ ($\approx 18.46$~МБ).
    \item \textbf{Эффективность BatchNormalization:} Экспериментально доказано, что при SGD с моментом ($\eta = 0.015$, $\beta = 0.9$) пакетная нормализация устраняет ковариационный сдвиг: модель с нормализацией достигает точности Top-5 свыше 89\,\% к 5-й эпохе, тогда как без нее сеть стагнирует на уровне менее 1.5\,\%.
    \item \textbf{Оценка на генеральной выборке:} На тестовом множестве из 6149 изображений достигнуты метрики: $\text{Top-1} = 27.87\,\%$, $\text{Top-5} = 55.07\,\%$ ($\text{Test Loss} = 3.1244$), что превосходит случайный базис в 28.4 и 11.2 раза соответственно при обучении всего на 10 примерах на класс.
    \item \textbf{Модуль инференса:} Реализован алгоритм классификации произвольных изображений цветов с построением гистограммы вероятностей наиболее правдоподобных классов Top-5.
\end{enumerate}

Разработанный комплекс подтвердил высокую эффективность сверточных архитектур с пакетной нормализацией в задачах компьютерного зрения со сложной межклассовой вариативностью.

